\documentclass[11pt]{article}

\usepackage[margin=1in]{geometry}
\usepackage{amsmath,amssymb,mathtools}
\usepackage{booktabs}
\usepackage{tabularx}
\usepackage{enumitem}
\usepackage{xcolor}
\usepackage{microtype}
\usepackage{url}
\usepackage[hidelinks]{hyperref}

\providecommand{\reportbuilddatetime}{unknown}
\InputIfFileExists{lps_coupled_support_chart_dimension_selection_report_build_info.tex}{}{}

\newcommand{\R}{\mathbb{R}}
\newcommand{\Gfull}{\mathcal G_{\mathrm{full}}}
\newcommand{\Gsparse}{\mathcal G_{\mathrm{sparse}}}
\newcommand{\argmin}{\operatorname*{arg\,min}}
\newcommand{\CV}{\operatorname{CV}}
\newcommand{\RMSE}{\operatorname{RMSE}}
\newcommand{\TruthRMSE}{\operatorname{TruthRMSE}}
\newcommand{\code}[1]{\texttt{#1}}

\title{\textbf{Coupled Support-Size and Chart-Dimension Selection in LPS}\\[3pt]
\large Definitions of \code{auto}, \code{local.auto}, \code{full\_kd}, and
\code{sparse\_kd}, with CSD5--CSD9 results}
\author{\small Source:\\[-1pt]
\begin{minipage}{0.92\textwidth}
\centering\footnotesize\ttfamily
\detokenize{~/current_projects/geosmooth/dev/methods/lps/status/}\\
\detokenize{lps_coupled_support_chart_dimension_selection_report.tex}
\end{minipage}}
\date{Build \reportbuilddatetime}

\begin{document}
\maketitle

\begin{abstract}
This report records the current meaning of the LPS support-size and chart-dimension
selection policies that appear in the CSD experiments.  Its main purpose is to
separate two ideas that are easy to confuse.  The existing \code{local.auto}
policy estimates an anchor-specific chart-dimension field, so different anchors
may use different local dimensions.  The current \code{full\_kd} and
\code{sparse\_kd} policies do not do this.  They select one global pair
$(k^\star,d^\star)$ for a fitted task, where $k$ is support size and $d$ is chart
dimension.  The CSD5--CSD9 experiments show that coupled global selection over
$(k,d)$ can improve over the older one-axis policies, but also that exact
inner-CV minimization over the full Cartesian grid is not a solved selector.
The strongest current evidence supports a prospective validation of a modest
near-tie rule, not an immediate default change.
\end{abstract}

\section{Purpose and executive summary}

The local polynomial smoother (LPS) fits a local polynomial in a local coordinate
chart around each prediction point.  Two tuning choices control the geometry of
each chart.  The support size $k$ controls the scale of the neighborhood, and the
chart dimension $d$ controls how many local principal directions are kept.  These
two quantities are statistically coupled: a larger neighborhood can support a
larger local dimension, while a small neighborhood often requires a smaller
dimension for numerical stability and for avoiding variance inflation.

The CSD series, short for coupled support-size and chart-dimension, was created
to study this coupled selection problem.  The central distinction is the
following.
\[
  \text{\code{local.auto}:}\qquad d_a \text{ may vary with anchor } a,
\]
whereas
\[
  \text{\code{full\_kd} and \code{sparse\_kd}:}\qquad
  k_a=k^\star,\quad d_a=d^\star \quad \text{for all anchors } a.
\]
Thus \code{full\_kd} is not an anchor-specific automatic dimension method.  It is
a full Cartesian inner-CV selector over global numeric pairs $(k,d)$.

The empirical lesson so far is mixed but useful.  On the CSD6 expanded suite,
\code{full\_kd} had lower median RMSE ratio than \code{auto} and \code{local.auto}
but remained meaningfully worse than the synthetic-reference candidate that
minimizes known-truth outer RMSE over the full numeric $(k,d)$ grid.  This
reference is defined below as the truth-facing full-grid oracle.  The
offline CSD9 replay suggested that a conservative one-percent near-tie rule
favoring lower dimension and non-boundary candidates slightly improved over exact
CV-minimum selection.  The improvement was modest: the median truth ratio moved
from $1.271$ to $1.252$, and the count of tasks with selected truth ratio above
$1.5$ moved from $16$ to $15$.  This is enough to motivate a prospective
validation run, but not enough to make the rule a package default.

\section{Local charts and the basic LPS object}

Let $X=\{x_1,\ldots,x_n\}\subset\R^p$ be the observed feature cloud, and let
$y_i$ be the response associated with $x_i$.  For an anchor $a$, the support
neighborhood at support size $k$ is
\[
  U_a(k)=\{x_a\}\cup\{\text{the } k-1 \text{ nearest neighbors of } x_a\}.
\]
In the current CSD experiments, this is a nearest-neighbor support in the input
geometry used by the LPS call.  For a local-PCA chart, the points in $U_a(k)$ are
centered at $x_a$ and projected onto their first $d$ principal directions.  The
resulting local coordinates are denoted
\[
  z_{a,i}^{(d)}\in\R^d,\qquad x_i\in U_a(k).
\]
LPS then builds a local polynomial design in these coordinates.  If the local
polynomial degree is $g$, the full monomial feature count, including the
intercept, is
\[
  q(d,g)=\binom{d+g}{g}.
\]
For example, if $g=1$, then $q(d,1)=d+1$.  If $g=2$, then $q(d,2)$ includes the
intercept, all linear terms, all squared terms, and all cross terms.

The local fit at anchor $a$ is a weighted local regression problem.  Its exact
numerical backend can be monomial, weighted-QR, or orthogonal-polynomial
stabilized, but conceptually it estimates local coefficients $\widehat\beta_a$
from weighted observations in $U_a(k)$ and returns the fitted value at the anchor.
With centered coordinates, the anchor prediction is the fitted intercept.

\section{Feasible numeric pairs}

A numeric candidate $(k,d)$ must leave enough observations to fit the local
design.  CSD uses the conservative prefit feasibility rule
\[
  k \ge q(d,g)+m,
\]
where $m\ge 0$ is a design margin.  In the CSD5--CSD9 experiments,
\[
  g=1,\qquad m=2,\qquad k\in\{15,\ldots,35\},\qquad d\in\{1,\ldots,8\}.
\]
Since $q(d,1)=d+1$, every planned pair in this grid is feasible:
\[
  d+1+2\le 11 \le 15
  \qquad\text{for all } d\le 8.
\]
This rule is intentionally only a prefit screen.  A candidate can still fail
inside the numerical fit if the realized weighted design is rank-deficient or
ill-conditioned.

\section{Candidate selection and synthetic evaluation}

This section defines the selection objects used throughout the report.  It is
included because the CSD labels are otherwise too easy to read as implementation
shorthand.

An \emph{outer task} is one fitted comparison unit.  In CSD5--CSD9, an outer
task is determined by a synthetic dataset, a random repetition, and an outer
cross-validation fold.  Let $T_s$ denote the outer-training indices for task
$s$, and let $H_s$ denote the outer-heldout indices.  Candidate selection is
performed using only the observed responses on $T_s$.  The known synthetic truth
$f_i$ is used only later, for auditing.

Inside $T_s$, the LPS fitting code creates inner validation folds.  For an inner
fold $\ell$, let $V_{s\ell}\subset T_s$ be the validation indices and
$T_{s\ell}=T_s\setminus V_{s\ell}$ be the corresponding inner-training indices.
For a candidate $\theta$, LPS is fit on $T_{s\ell}$ and predicts the held-out
responses at $V_{s\ell}$.  The inner-CV score is
\[
  C_s(\theta)
  =
  \left\{
  \frac{\sum_{\ell}\sum_{i\in V_{s\ell}}
  \left(y_i-\widehat f^{(-\ell)}_{s,\theta}(x_i)\right)^2}
  {\sum_{\ell}|V_{s\ell}|}
  \right\}^{1/2}.
\]
Here $\widehat f^{(-\ell)}_{s,\theta}$ is the LPS fit trained without the
responses in $V_{s\ell}$.  The selected candidate is
\[
  \widehat\theta_s
  =
  \argmin_{\theta\in\mathcal A_s} C_s(\theta),
\]
where $\mathcal A_s$ is the candidate set for the selector being used.  If two
candidates are exactly tied, the implementation uses deterministic candidate
ordering.  In practice exact ties are uncommon except in deliberately simplified
smoke cases.

The meaning of $\theta$ depends on the selector.  For \code{auto} and
\code{local.auto}, the candidate is the support size $k$, because chart
dimensions are resolved by the automatic dimension rule.  Thus
\[
  \mathcal A_s^{\mathrm{auto}}=\{15,16,\ldots,35\},
  \qquad
  \theta=k.
\]
For \code{full\_kd} and \code{sparse\_kd}, the candidate is a numeric pair
\[
  \theta=(k,d),
\]
where $k$ is support size and $d$ is a scalar chart dimension used by all
anchors in that fit.

After $\widehat\theta_s$ is selected, LPS is refit on all of $T_s$ using that
candidate.  The fitted model is then evaluated on $H_s$.  In real applications,
one would report an outer observed prediction score against $y_i$.  In the CSD
synthetic experiments, the main audit score is truth RMSE against the known
synthetic function:
\[
  R_{s,m}
  =
  \left\{
  \frac{1}{|H_s|}
  \sum_{i\in H_s}
  \left(\widehat f_{s,m}(x_i)-f_i\right)^2
  \right\}^{1/2},
\]
where $m$ is the selector label and $\widehat f_{s,m}$ is the final fit selected
by that method.

The \emph{truth-facing full-grid oracle} is a synthetic-only reference.  For
each outer task $s$, we imagine fitting every feasible numeric pair
$(k,d)\in\Gfull$, evaluating each one against the known truth on $H_s$, and
keeping the best truth RMSE:
\[
  R_s^\star
  =
  \min_{(k,d)\in\Gfull}
  \left\{
  \frac{1}{|H_s|}
  \sum_{i\in H_s}
  \left(\widehat f_{s,k,d}(x_i)-f_i\right)^2
  \right\}^{1/2}.
\]
It is an \emph{oracle} because it uses $f_i$, which is unavailable in real data.
It is \emph{truth-facing} because it minimizes truth error rather than inner CV
error.  It is \emph{full-grid} because the minimization is over the complete
numeric grid $\Gfull$, not over a sparse candidate subset.  A deployable
selector such as \code{full\_kd} may search the same grid, but it must select by
$C_s(k,d)$ rather than by $R_s(k,d)$.

The main relative accuracy diagnostic is
\[
  \kappa_{s,m}
  =
  \frac{R_{s,m}}{R_s^\star+\epsilon},
  \qquad
  \epsilon=10^{-12}.
\]
A value $\kappa_{s,m}=1$ means that method $m$ matched the truth-facing
full-grid oracle on that outer task.  A value $1.25$ means the selected fit has
25 percent larger truth RMSE than the oracle.  A value $2$ means twice the
oracle truth RMSE.

\section{Automatic chart-dimension rules}

The automatic dimension policies are based only on the observed feature cloud
$X$, the support size $k$, and the local polynomial degree $g$.  They do not use
the response $y$ and do not use the synthetic truth $f$.

At an anchor $a$, form the local-PCA singular values of the centered support
$U_a(k)$:
\[
  \sigma_{a1}\ge \sigma_{a2}\ge\cdots\ge 0.
\]
Let $e_{ar}=\sigma_{ar}^2$ be the local spectral energy.  The implementation
computes two candidate dimensions.  The cumulative-energy dimension is the first
$r$ satisfying
\[
  \frac{\sum_{j=1}^{r} e_{aj}}{\sum_{j\ge 1} e_{aj}}
  \ge \tau_{\mathrm{var}},
\]
with default variance threshold $\tau_{\mathrm{var}}=0.95$.  The eigengap
dimension is the index $r$ at which the ratio
\[
  \frac{\sigma_{ar}}{\max(\sigma_{a,r+1},\epsilon_{\mathrm{sv}})}
\]
is largest, provided that the largest ratio is at least the default eigengap
threshold $\tau_{\mathrm{gap}}=4$.  If the eigengap rule does not find a
sufficiently large gap, it is ignored.  The anchor-level automatic dimension is
then
\[
  \widehat d_a(k)
  =
  \begin{cases}
  \min\{d^{\mathrm{var}}_a(k),d^{\mathrm{gap}}_a(k)\},
    &\text{if the eigengap dimension is available},\\
  d^{\mathrm{var}}_a(k),
    &\text{otherwise}.
  \end{cases}
\]
The result is clipped by the local design feasibility cap.  For a support with
$k$ points, the cap is the largest integer $r$ such that
\[
  q(r,g)=\binom{r+g}{g}\le k-1.
\]
The subtraction by one is conservative: it leaves one degree of freedom beyond
the raw number of local polynomial columns.  This cap prevents the automatic
dimension rule from requesting more local polynomial columns than the support
can stably fit.

The \code{local.auto} policy uses this anchor-level value directly:
\[
  d_a=\widehat d_a(k).
\]
The \code{auto} policy pools the same anchor-level diagnostics into one global
dimension.  In the current implementation, the global automatic dimension is the
median of the successful anchor-level values, again clipped to the feasible
range.  For speed, the global \code{auto} helper uses at most 60 representative
anchors when $n$ is larger than 60, while \code{local.auto} evaluates all
anchors because it needs a dimension value for each one.  The pooled global
dimension is
\[
  \widehat d_{\mathrm{auto}}(k)
  =
  \operatorname{median}_{a:\widehat d_a(k)\ \mathrm{available}}
  \widehat d_a(k),
\]
after integer conversion and feasibility clipping.  If no anchor-level estimate
is available, the code uses a conservative fallback dimension allowed by the
support-size and degree constraints.

\section{The four selector families}

This section defines the four policy labels used in the CSD reports.  The most
important axis is whether the policy chooses a global dimension, an
anchor-specific dimension field, or a global numeric pair.

\subsection{\code{auto}: one global automatic chart dimension}

The \code{auto} policy asks the package to estimate a single chart dimension from
the observed input geometry.  Because the automatic dimension depends on the
support size, write this value as
\[
  \widehat d_{\mathrm{auto}}(k).
\]
Once it is resolved, the same chart dimension is used at every anchor:
\[
  d_a=\widehat d_{\mathrm{auto}}(k)\qquad\text{for all anchors }a.
\]
The support-size candidate set in CSD5--CSD9 is
\[
  \mathcal K=\{15,16,\ldots,35\}.
\]
For each $k\in\mathcal K$, the code resolves
$\widehat d_{\mathrm{auto}}(k)$ from the local-PCA spectra of $X$, fits LPS in
inner cross-validation using the global dimension
$d_a=\widehat d_{\mathrm{auto}}(k)$, and computes the inner-CV score
\[
  C_s^{\mathrm{auto}}(k)
  =
  C_s\!\left(k,\widehat d_{\mathrm{auto}}(k)\right).
\]
The selected support size is
\[
  \widehat k_s^{\mathrm{auto}}
  =
  \argmin_{k\in\mathcal K} C_s^{\mathrm{auto}}(k).
\]
The final outer-training fit then uses
\[
  k_a=\widehat k_s^{\mathrm{auto}},
  \qquad
  d_a=\widehat d_{\mathrm{auto}}\!\left(\widehat k_s^{\mathrm{auto}}\right)
  \qquad\text{for all anchors }a.
\]
Thus, in the CSD comparison, \code{auto} is a one-dimensional CV search over
support size, with a geometry-derived global chart dimension attached to each
candidate support size.

This policy is deployable and simple, but it assumes that a single working local
dimension is adequate for the whole dataset.  That is plausible for homogeneous
manifold-like data and less plausible for stratified or non-manifold data, such
as simplex-boundary geometries or unions of pieces with different intrinsic
dimensions.

\subsection{\code{local.auto}: anchor-specific automatic chart dimensions}

The \code{local.auto} policy estimates chart dimension separately by anchor:
\[
  d_a=\widehat d_a(k).
\]
The same local spectral rule from the previous section is applied at each anchor
and at each candidate support size.  For a fixed $k$, the local polynomial model
at anchor $a$ is fit in $\R^{\widehat d_a(k)}$.

The support size is not anchor-specific in the current CSD setup.  The candidate
set is again $\mathcal K=\{15,\ldots,35\}$.  For each $k\in\mathcal K$, the
inner-CV fit uses
\[
  k_a=k,\qquad d_a=\widehat d_a(k),
\]
and computes a score denoted $C_s^{\mathrm{local.auto}}(k)$.  The selected
support size is
\[
  \widehat k_s^{\mathrm{local.auto}}
  =
  \argmin_{k\in\mathcal K} C_s^{\mathrm{local.auto}}(k).
\]
The final fitted model uses one selected support size for all anchors but keeps
anchor-specific chart dimensions:
\[
  k_a=\widehat k_s^{\mathrm{local.auto}},
  \qquad
  d_a=\widehat d_a\!\left(\widehat k_s^{\mathrm{local.auto}}\right).
\]

This policy is the one that should be described as anchor-specific.  It can adapt
to geometries where some regions look one-dimensional and other regions look
two- or three-dimensional.  Its cost is that the local designs have varying
numbers of columns, and its statistical risk is that noisy local dimension
estimates can create unstable chart fields.

\subsection{\code{full\_kd}: full global Cartesian selector over numeric pairs}

The \code{full\_kd} policy defines a full Cartesian candidate set
\[
  \Gfull
  =
  \{15,16,\ldots,35\}\times\{1,2,\ldots,8\},
\]
after feasibility filtering.  For each candidate $(k,d)\in\Gfull$, LPS is fit
with
\[
  k_a=k,\qquad d_a=d\qquad\text{for all anchors }a.
\]
The inner-CV score for a candidate is denoted
\[
  C_s(k,d),
\]
where $s$ indexes the outer task, meaning a particular dataset, repetition, and
outer fold.  To compute this number, the code repeats the same fold procedure
described above: fit on each inner-training subset, predict the corresponding
inner-validation subset, square the validation residuals, and average them over
all inner validation points.  The current exact CV-minimum version of
\code{full\_kd} selects
\[
  (\widehat k_s,\widehat d_s)
  =
  \argmin_{(k,d)\in\Gfull} C_s(k,d).
\]
The selected pair is global for task $s$.  It is not a vector
$\{(k_a,d_a):a=1,\ldots,n\}$.  Once selected, all anchors in that fitted model
use the same support size and the same numeric chart dimension.

In the CSD experiments, \code{full\_kd} has two roles.  First, it is a
computationally expensive but auditable candidate selector using inner CV.
Second, when scored against known synthetic truth across all pairs, it provides a
truth-facing reference grid.  These roles must not be confused.  The deployable
selector can only use $C_s(k,d)$; the truth-facing oracle uses
$R_s(k,d)$ and is available only in synthetic experiments.

\subsection{\code{sparse\_kd}: sparse global numeric selector}

The \code{sparse\_kd} policy is a cheaper approximation to \code{full\_kd}.  It
uses a smaller candidate set
\[
  \Gsparse\subset\Gfull.
\]
In the CSD5 and CSD6 runs, the sparse grid was constructed mechanically from the
same planned grids.  The support skeleton retained the minimum, middle, and
maximum planned support sizes:
\[
  \mathcal K_{\mathrm{sparse}}=\{15,25,35\}.
\]
The dimension skeleton retained the low dimensions $1$ and $2$ and the largest
planned dimension available after feasibility clipping:
\[
  \mathcal D_{\mathrm{sparse}}=\{1,2,8\}.
\]
Their Cartesian product produced
\[
  \Gsparse
  =
  \mathcal K_{\mathrm{sparse}}\times \mathcal D_{\mathrm{sparse}},
\]
so 9 candidates were evaluated rather than the 168 candidates in
$\Gfull=\{15,\ldots,35\}\times\{1,\ldots,8\}$.  As with \code{full\_kd}, a
selected sparse candidate is a single global pair:
\[
  (\widehat k_s,\widehat d_s)\in\Gsparse.
\]
Every anchor in the fitted model then uses this same selected pair.

The engineering motivation for \code{sparse\_kd} is PCA-coordinate reuse.  For a
fixed support size and kernel, the package can build local PCA coordinates once
at the largest requested numeric dimension and reuse leading columns for smaller
numeric dimensions.  This makes numeric chart-dimension grids far cheaper than
rebuilding every chart from scratch.  The statistical risk is that a sparse grid
can miss the good region of the full grid, so it must be benchmarked against
\code{full\_kd} before being treated as a routine selector.

\section{How to read the CSD tables}

The CSD tables below use the notation just defined.  The column called RMSE
ratio is the median of $\kappa_{s,m}$ over matched outer tasks.  The column
called relative regret is
\[
  100\,(\kappa_{s,m}-1).
\]
Thus, relative regret is measured in percent above the truth-facing full-grid
oracle.  A relative regret of 0 percent means the selected method matched the
oracle on that task; 25 percent means the selected truth RMSE was 1.25 times the
oracle truth RMSE.

\section{What the CSD5--CSD9 experiments show}

The CSD experiments are not one experiment repeated under different names.  They
answer a sequence of increasingly specific questions.

\subsection{CSD5: first coupled-grid smoke}

CSD5 compared \code{auto}, \code{local.auto}, \code{sparse\_kd}, and
\code{full\_kd} on a compact smoke suite.  Table~\ref{tab:csd5} gives the median
strategy summaries.  The main signal was that both numeric coupled selectors did
better than the one-axis automatic policies in this small run, and
\code{sparse\_kd} was surprisingly competitive while using far fewer candidates.

\begin{table}[htbp]
\centering
\caption{CSD5 smoke summary.  The RMSE ratio is the median selected truth RMSE
divided by the truth-facing full-grid oracle on matched outer tasks.  Candidate
and PCA counts are medians per outer fit.}
\label{tab:csd5}
\begin{tabular}{lrrrr}
\toprule
Strategy & RMSE ratio & Relative regret (\%) & Candidates & PCA builds\\
\midrule
\code{auto}       & 2.022 & 102.2 & 21  & 21\\
\code{local.auto} & 2.114 & 111.4 & 21  & 21\\
\code{sparse\_kd} & 1.242 &  24.2 &  9  &  3\\
\code{full\_kd}   & 1.334 &  33.4 & 168 & 21\\
\bottomrule
\end{tabular}
\end{table}

Because CSD5 was small, its main value was not the precise ranking.  Its value
was showing that coupled $(k,d)$ search was worth testing more seriously and
that PCA reuse made sparse numeric grids cheap enough to consider.

\subsection{CSD6: expanded relative-regret evaluation}

CSD6 expanded the geometry suite to 8 datasets, 2 repetitions, and 3 outer folds,
for 48 matched outer tasks per strategy.  Table~\ref{tab:csd6} summarizes the
median results.  In this broader run, \code{full\_kd} had the lowest median RMSE
ratio, while \code{sparse\_kd} remained close and much cheaper.

\begin{table}[htbp]
\centering
\caption{CSD6 expanded summary.  The full Cartesian grid evaluated 168 numeric
$(k,d)$ candidates per task, while the sparse grid evaluated 9.  The lower RMSE
ratio for \code{full\_kd} indicates better median accuracy, but the ratio remains
well above 1, so exact CV minimization over the full grid is not oracle-like.}
\label{tab:csd6}
\begin{tabular}{lrrrr}
\toprule
Strategy & RMSE ratio & Relative regret (\%) & Candidates & PCA builds\\
\midrule
\code{auto}       & 1.374 & 37.4 & 21  & 21\\
\code{local.auto} & 1.410 & 41.0 & 21  & 21\\
\code{sparse\_kd} & 1.298 & 29.8 &  9  &  3\\
\code{full\_kd}   & 1.271 & 27.1 & 168 & 21\\
\bottomrule
\end{tabular}
\end{table}

The CSD6 conclusion is therefore not that \code{full\_kd} solves selection.  It
is that global coupled numeric selection improves the median picture relative to
the older automatic one-axis policies, but still leaves substantial selection
regret.

\subsection{CSD7: diagnosing failures}

CSD7 asked why the CSD6 selectors still missed.  The failure diagnostics showed
that poor outcomes were not only sparse-grid misses.  Some failures occurred even
when \code{full\_kd} was available.  In particular, exact CV-minimum selection
sometimes chose a candidate far from the truth-facing best candidate even though
the better candidate was present in the full grid.

This changed the next question.  Before CSD7, it was tempting to ask whether the
sparse grid was too sparse.  After CSD7, the more important question became
whether the inner-CV surface itself was selecting too aggressively among near-tied
candidates or boundary candidates.

\subsection{CSD8: candidate-level CV surface audit}

CSD8 persisted the full candidate-level inner-CV surface and joined it to the
CSD6 truth-facing grid.  Each outer task had 168 candidate rows, giving 8064
candidate rows in total.  This made it possible to compare
\[
  C_s(k,d)
  \qquad\text{with}\qquad
  R_s(k,d)
\]
candidate by candidate.

The key finding was that low CV score did not always imply low truth RMSE.  Some
tasks had several candidates with nearly indistinguishable CV scores, while their
truth RMSE values differed substantially.  This supported the idea that exact
CV-minimum selection may be too brittle, especially when the winning candidate
lies on the edge of the grid or when many candidates are essentially tied in CV.

\subsection{CSD9: offline robust replay rules}

CSD9 replayed simple deterministic policies on the saved CSD8 surfaces.  No
models were refit.  The baseline was exact CV minimum:
\[
  (\widehat k_s,\widehat d_s)
  =
  \argmin_{(k,d)\in\Gfull} C_s(k,d).
\]
The most successful replay rule was the one-percent low-dimension rule.  It
first formed the near-tie set
\[
  T_s(0.01)
  =
  \{(k,d)\in\Gfull: C_s(k,d)\le 1.01 \min_{(k',d')\in\Gfull} C_s(k',d')\}.
\]
If non-boundary candidates existed inside this set, boundary candidates were
discarded.  Among the remaining candidates, the rule preferred smaller $d$, then
$k$ closest to the middle of the support grid, then smaller CV score.

\begin{table}[htbp]
\centering
\caption{CSD9 policy replay summary.  The best replay was the one-percent
low-dimension rule, but the gain over exact CV minimum was modest.  A severe
miss means selected truth ratio above 1.5.}
\label{tab:csd9}
\begin{tabular}{lrrrr}
\toprule
Policy & Median truth ratio & Severe misses & Boundary rate & Median $k,d$\\
\midrule
1\% low-d      & 1.252 & 15 & 0.417 & $(19.5,6)$\\
CV minimum     & 1.271 & 16 & 0.521 & $(19,6)$\\
low-d penalty  & 1.284 & 17 & 0.271 & $(19,5.5)$\\
boundary penalty & 1.284 & 17 & 0.292 & $(20,6)$\\
3\% large-k    & 1.328 & 17 & 0.271 & $(21.5,6)$\\
3\% low-d      & 1.336 & 18 & 0.271 & $(19,5)$\\
5\% low-d      & 1.359 & 17 & 0.271 & $(21,5)$\\
\bottomrule
\end{tabular}
\end{table}

CSD9 therefore gives a useful but deliberately cautious recommendation.  The
one-percent low-dimension rule is worth testing prospectively.  It should not yet
replace the current selector by default, because it was chosen after seeing the
saved CSD8 surface and the paired median deltas were often exactly zero because
many replay rules selected the same candidate as exact CV minimum.

\section{Current interpretation}

The CSD work clarifies the current taxonomy of dimension and support policies.
\code{auto} is a global automatic dimension policy.  \code{local.auto} is an
anchor-specific automatic dimension policy.  \code{full\_kd} is a global
full-grid numeric selector over $(k,d)$.  \code{sparse\_kd} is a cheaper global
numeric selector over a sparse subset of that grid.

The phrase ``full grid'' should therefore be read carefully.  It means full in
the candidate-grid sense:
\[
  \Gfull=\{15,\ldots,35\}\times\{1,\ldots,8\}.
\]
It does not mean full in the sense of allowing every anchor to choose its own
$(k_a,d_a)$.  An anchor-specific coupled selector would be a different method:
\[
  a\mapsto (k_a,d_a).
\]
That method is conceptually attractive, but it is not what CSD5--CSD9 tested.

\section{Recommended next step}

The next useful step is a prospective CSD10 validation run.  It should compare
at least three policies on the same outer tasks:
\begin{align*}
  &\text{exact CV minimum over }\Gfull,\\
  &\text{the 1\% low-dimension near-tie replay rule made prospective},\\
  &\text{a sparse-grid approximation}.
\end{align*}
The one-percent rule must be implemented inside the selector, not replayed after
the fact.  The report should preserve the candidate-level inner-CV surface and
truth-facing outer scores so that the same diagnostic questions remain auditable.

A later, larger methodological direction is an anchor-specific coupled
scale-dimension field:
\[
  a\mapsto (k_a,d_a).
\]
That would be closer in spirit to \code{local.auto}, but it would require new
regularization or stabilization machinery.  It should not be described as the
current \code{full\_kd} method.

\section*{Artifact pointers}

The main artifacts summarized here are:
\begingroup
\footnotesize
\sloppy
\begin{itemize}[leftmargin=1.4em]
\item CSD plan:
\path{~/current_projects/geosmooth/dev/methods/lps/specs/coupled_support_chart_dimension_selection_plan_2026-07-08.md}.
\item CSD5 report:
\path{~/current_projects/geosmooth/dev/methods/lps/reports/csd5_coupled_kd_evaluation_20260708/csd5_coupled_kd_evaluation_report.html}.
\item CSD6 report:
\path{~/current_projects/geosmooth/dev/methods/lps/reports/csd6_expanded_relative_regret_20260708/csd6_expanded_relative_regret_report.html}.
\item CSD7 report:
\path{~/current_projects/geosmooth/dev/methods/lps/reports/csd7_task_failure_diagnostics_20260708/csd7_task_failure_diagnostics_report.html}.
\item CSD8 report:
\path{~/current_projects/geosmooth/dev/methods/lps/reports/csd8_candidate_cv_surface_audit_20260708/csd8_candidate_cv_surface_audit_report.html}.
\item CSD9 report:
\path{~/current_projects/geosmooth/dev/methods/lps/reports/csd9_robust_cv_selection_policy_audit_20260708/csd9_robust_cv_selection_policy_audit_report.html}.
\end{itemize}
\endgroup

\end{document}
